\begin{textsample}{POS Dim 2 – ba – Score 22.00 – ba/ba\_em\_10.txt}  \label{ex:f2_pos_040}
5.5.
TERRITORIALIDADE O Estado da Bahia possui ampla pluralidade cultural e de agrupamentos humanos dos quais emergem as expressões culturais das diversas nações africanas, dos povos indígenas, dos sertanejos, dos grapiúnas, da região do São Francisco, dos povos ciganos, dos quilombolas, dentre outros, com suas formas de vida autênticas, que revelam a capacidade de seus habitantes de expressarem suas identidades territoriais, por meio da sua história, culturas, saberes, conhecimentos, produções, dimensões fundamentais de sua identificação comunitária, dos modos de viver, conviver e produzir, das maneiras de ser, sentir e querer.
Essa pluralidade e diversidade cultural são expressas nos 417 \textbf{municípios} do Estado, organizados em 27 Territórios de Identidade, com dinâmicas demográficas específicas e aspectos diversificados quanto à sua configuração ambiental, geográfica e de desenvolvimento econômico e político que demandam \textbf{políticas} adequadas e singulares à composição demográfica e de identidade da população.
O conceito de Território de Identidade, previsto na Lei nº 13.214/2014, dá centralidade à \textbf{política} de desenvolvimento territorial do Estado da Bahia, a saber : > entende -se por Território de Identidade a unidade de planejamento de \textbf{políticas} públicas do Estado da Bahia, constituído por agrupamentos identitários \textbf{municipais}, geralmente contíguos, formado de acordo com critérios sociais, culturais, econômicos e geográficos, reconhecido pela sua população como o espaço historicamente construído ao qual pertencem, com identidade que amplia as possibilidades de coesão social e territorial ( BAHIA, 2014 ).
Essa organização política do Estado da Bahia, por Territórios de Identidade, busca estabelecer articulações entre o Governo Estadual e a sociedade, visando a promoção do desenvolvimento sustentável e a \textbf{equidade} no acesso às oportunidades sociais, econômicas e à cidadania.
Tal organização objetiva estimular a cooperação entre os atores e atrizes sociais e fortalecer as várias dinâmicas sociais e econômicas, articulando -as a partir de um mesmo referencial de planejamento de \textbf{políticas} de desenvolvimento e de iniciativas locais.
Segundo Milton Santos, > o território pode ser visto como uma materialidade ( configuração territorial ) cuja apreensão por meio dos sentidos caracteriza -o como paisagem.
Como uma forma política e econômica a caracterizar o espaço, categoria, objeto e totalidade social, o território contém subespaços que seriam as regiões ( SANTOS apud FUINI, 2015 ).
Considerando -se, ainda, a relação entre a Política Territorial e a Política Estadual de Convivência com o Semiárido ( Lei Estadual 13.572/2016 ), constata -se que o Governo adotou instrumentos que articulam as suas \textbf{políticas} públicas com os territórios, promovendo uma maior aderência às distintas necessidades territoriais e ampliando assim a efetividade das ações governamentais, além de promover o \textbf{fortalecimento} de uma democracia mais participativa.
Trata -se, portanto, do amadurecimento e execução de \textbf{políticas} públicas voltadas à convivência com o semiárido com o objetivo de institucionalizar uma \textbf{política} pública que estabelece um conjunto de princípios e \textbf{diretrizes} para o desenvolvimento da região semiárida.
É um instrumento de gestão e planejamento \textbf{intersetorial} e transversal de \textbf{políticas} e programas governamentais e ações da sociedade civil, cujo objetivo é fortalecer e promover a autonomia das populações, com a inclusão socioprodutiva e geração de emprego e renda.
É sabido que o semiárido ocupa 70 \% do território da Bahia, o que equivale a uma área de cerca de 390 mil quilômetros quadrados, com uma população de mais de 6,5 milhões de pessoas, distribuídas em aproximadamente 270 \textbf{municípios}.
Logo, mais do que necessário, torna -se imprescindível observar o Art. 4º do \textbf{Capítulo} II, que trata dos princípios da Política Estadual de Convivência com o Semiárido, a saber : I. garantia da permanência sustentável, digna e cidadã das populações em seus territórios ; II. \textbf{garantia} do acesso e da permanência na terra ; III. universalização do acesso à água ; IV. \textbf{garantia} dos usos múltiplos dos recursos hídricos de forma racional ; V. conservação e preservação da biodiversidade e uso sustentável dos recursos naturais ; VI. universalidade e \textbf{equidade} no acesso às \textbf{políticas} públicas que promovam a convivência com o semiárido ; VII. intersetorialidade e \textbf{transversalidade} das \textbf{políticas} públicas ; VIII. transparência e descentralização ; IX. participação e controle social ; X. valorização, respeito e proteção às diversidades social, cultural, ambiental, econômica, étnico-racial, geracional e de gênero ; XI. direito à educação pública, gratuita, integral, de \textbf{qualidade}, contextualizada ao semiárido ; XII. direito à saúde como suporte à \textbf{qualidade} de vida ; XIII. economia solidária ; XIV. agricultura sustentável e agroecológica ; XV. produção associada ao turismo.
Em consonância com a Política Territorial e a Política de Convivência com o semiárido, a Secretaria da Educação do Estado da Bahia \textbf{instituiu} os 27 Núcleos Territoriais de Educação ( NTE ), a partir do Decreto nº 17.377/2017, considerando as diferentes dimensões e características dos Territórios de Identidade, agregando, assim, sentimento de pertença da população aos locais onde vivem, sendo este um direcionamento fundamental para a adequação das \textbf{políticas} educacionais \textbf{vigentes}, e que ganha destaque e centralidade na \textbf{política} curricular do Estado, por meio do Documento Curricular Referencial da Bahia ( DCRB ), disponibilizados em três volumes : - O Volume 1, da Educação Infantil e Ensino Fundamental, \textbf{homologado} em 19 de dezembro de 2019. - O Volume 2, \textbf{destinado} às \textbf{ofertas} do Ensino Médio Regular, da modalidade da Educação Profissional e Tecnológica e da \textbf{oferta} da Educação Integral, \textbf{homologado} em 29 de março de 2022. - O Volume 3, voltado para as etapas da Educação Infantil, Ensino Fundamental e Ensino Médio das Modalidades de Ensino.
Em fase de elaboração.
O Volume 1 do Documento Curricular Referencial da Bahia, para a Educação Infantil e Ensino Fundamental, já anunciava a territorialidade como centralidade para a organização curricular do Estado e organização curricular das Unidades Escolares, por meio dos seus Projetos Políticos Pedagógicos ( PPPs ), valorizando, evidenciando e contextualizando os saberes, conhecimentos da população local, suas lideranças e mestres/as Griôs², produções culturais, aspectos geográficos, históricos, biológicos, arranjos produtivos locais e a dinâmica da economia criativa, como aspectos pujantes e oportunos para a promoção de aprendizagens e para a produção de conhecimento local.
Segundo Farias, 2019 : > O entorno da unidade escolar, as bacias hidrográficas, as áreas naturais mais próximas, os aspectos da urbanização e da ação antrópica que modificam a paisagem natural, os grupos culturais locais, as associações, os pontos de encontro da comunidade – todos esses aspectos devem ser inseridos nas práticas \textbf{formativas} escolares, direcionando olhares, investigações, sem prejuízo do rigor científico, ao contrário, conduzindo o aprimoramento da pesquisa científica a partir de elementos da territorialidade.
Essa contextualização do território, dada como possibilidade a partir do desenvolvimento do Documento Curricular Referencial da Bahia, para a Educação Infantil e Ensino Fundamental e para o Ensino Médio e suas Modalidades, nas escolas, torna -se importante elemento crítico para a compreensão das identidades territoriais e para a observação dos fenômenos dinâmicos que promovem desterritorializações e reterritorializações de sujeitos, despertando o olhar para os conflitos subjacentes nas realidades em movimento, para o devir, para as “ sangrias ” daquilo que é conservador ou impedidor de \textbf{mobilidades} – subjetivas e concretas. ( BAHIA, 2019, p. 27 – grifos nossos ).
Destarte, compreender a territorialidade baiana, como um dos eixos \textbf{estruturantes} do DCRB, que visa impulsionar a organização dos \textbf{currículos} escolares do Estado, a partir das possibilidades e riquezas que existem nos Territórios de Identidade como importante elemento contextual para a formação básica dos/as estudantes baianos/as.

% matched lemmas: capítulo, currículo, destinar, diretriz, equidade, estruturante, formativo, fortalecimento, garantia, homologar, instituir, intersetorial, mobilidade, municipal, município, oferta, política, qualidade, transversalidade, vigente
\end{textsample}
